\noindent Кластерное разложение C46 с классической статистикой колебаний даёт
\textbf{14,90 ат.\% Al в \(\alpha\)} и
\textbf{21,92 ат.\% Al в \(\alpha_2\)}.
На 5120 узлах этот результат получен как на 13 окнах оригинала, так и на
уточнённой сетке из 25 окон. Условные 95\%-е интервалы на густой сетке
составляют \(14{,}861\)--\(14{,}933\) и \(21{,}867\)--\(21{,}978\)
ат.\% Al. Классический и квантовый варианты основаны на
одних декорациях C46 и статических энергиях потенциала Qi и соавторов.
В обоих расчётах объём флуктуирует, а \(c/a=1{,}56374\) фиксировано.

\subsection*{Сдвиг обеих границ при замене статистики}
На исходных 13 окнах квантовая модель даёт \(14{,}943/21{,}931\) ат.\%
Al. Архивные значения C46, \(14{,}958/21{,}958\), входят в интервалы
нового квантового расчёта. Переход к классике сдвигает границу
\(\alpha\) к меньшим содержаниям Al. Для \(\alpha_2\) интервал парной разности
включает нуль. Поле сосуществования уменьшается на \(2{,}015\)~мэВ на
замену Ti на Al. В таблице сдвиг отсчитан от квантовой границы, найденной
на том же наборе траекторий.

% RESULTS_TABLE

При переходе от 13 к 25 окнам абсолютные классические границы изменились
меньше чем на \(0{,}002\) п.п. На 25 окнах замена интерполяции PCHIP
линейными отрезками даёт парные сдвиги \(-0{,}036\) и \(-0{,}003\) п.п.
При уменьшении ячейки до 768 узлов сдвиг границы \(\alpha\) сохраняет
знак и возрастает по модулю. Для \(\alpha_2\) при обоих размерах
парный сдвиг сопоставим со статистической погрешностью.

\subsection*{Сравнение с прямым классическим MTP}
При тех же \(T,p,c/a\) расчёт с явными смещениями на 1400 атомах при
\(W<0{,}40\) дал \(16{,}083\) и \(22{,}159\) ат.\% Al. Относительно
классического C46 на авторской сетке 13 окон расхождение
«прямой MTP минус C46» составляет \(+1{,}187\) п.п. для \(\alpha\)
и \(+0{,}237\) п.п. для \(\alpha_2\). Для уточнённой сетки получены
\(+1{,}187\;[1{,}123;\,1{,}257]\) и \(+0{,}238\;[0{,}131;\,0{,}336]\)
п.п. Интервалы учитывают статистику двух независимых расчётов.
Поправка за статистику ядер слишком мала, чтобы объяснить расхождение
на Ti-богатой границе. В остаток входят усечение кластерного разложения,
квазигармоническое приближение, ограничения на смещения прямого MTP,
конечные размеры и способ извлечения составов.
